% Copyright (c) 2026 Geovana Neves. Licensed under CC BY 4.0 (https://creativecommons.org/licenses/by/4.0/). See LICENSE-AND-AUTHORSHIP.md.
%
\section{What the pproc asks the solver for}

\begin{frame}{Force plots: \code{[plots] parameters}}
\relax
{\ftstablesize
\begin{tabular}{@{}p{0.16\linewidth}p{0.25\linewidth}p{0.29\linewidth}p{0.208\linewidth}@{}}
\toprule
\textbf{In the pproc} & \textbf{The solver's PARAMETER} & \textbf{Units the package requests} & \textbf{Plots column} \\
\midrule
\code{CL}  & \code{CL}       & \code{COEFFICIENTS} & \code{CL\_<group>} \\
\code{CDI} & \code{CDI}      & \code{COEFFICIENTS} & \code{CDI\_<group>} \\
\code{CDO} & \code{CDO}      & \code{COEFFICIENTS} & \code{CDO\_<group>} \\
\code{CD}  & \code{CD}       & \code{COEFFICIENTS} & \code{CD\_<group>} \\
\code{FX}  & \code{FORCE\_X}  & \code{NEWTONS} & \code{FX\_<group>} \\
\code{FY}  & \code{FORCE\_Y}  & \code{NEWTONS} & \code{FY\_<group>} \\
\code{FZ}  & \code{FORCE\_Z}  & \code{NEWTONS} & \code{FZ\_<group>} \\
\code{MX}  & \code{MOMENT\_X} & \code{NEWTONS} & \code{MX\_<group>} \\
\code{MY}  & \code{MOMENT\_Y} & \code{NEWTONS} & \code{MY\_<group>} \\
\code{MZ}  & \code{MOMENT\_Z} & \code{NEWTONS} & \code{MZ\_<group>} \\
\bottomrule
\end{tabular}}
\begin{itemize}\scriptsize
  \item Write the short name. The pproc refuses \code{FORCE\_X}.
        Omitted \code{parameters} means all ten. There is no \code{CY}, \code{CM} or \code{CN} plot.
  \item The solver also offers \code{KILO-NEWTONS}, \code{POUND-FORCE} and
        \code{KILOGRAM-FORCE}. The package fixes one unit token per parameter.
        The plots table rescales coefficient columns by $(\mathrm{VREF}/\mathrm{VINF})^2$.
  \item Global \code{MRP} loads \code{FX .. MZ} feed the unsteady polar and rotor table.
        With \code{MRP} available, the run adds \code{MRP\_TOTAL} if no MRP group plots all six,
        and \code{ROTOR\_<ALIAS>} if none plots all six over that rotor's families.
\end{itemize}
{\tiny\color{SoftGray} Sources under \code{src/pyflightstream/}:
\code{cases/\_\_init\_\_.py}, \code{FORCE\_PLOT\_PARAMETERS}, \code{PlotsSpec};
\code{commands/unsteady\_solver.yaml}, \code{UNSTEADY\_SOLVER\_NEW\_FORCE\_PLOT};
\code{cases/workflows.py}, \code{\_pproc\_plots}, \code{\_plot\_each\_rotors\_own\_history};
\code{post/products.py}, \code{write\_plots\_table}.\par}
\end{frame}

\begin{frame}{Probes: \code{[[probes]] parameters}}
\small The pproc and the solver use the same spelling:
\par\vspace{2mm}
{\ftstablesize
\begin{tabular}{@{}llll@{}}
\code{CP\_FREE} & \code{CP\_REF} & \code{MACH} & \code{VELOCITY} \\
\code{VX} & \code{VY} & \code{VZ} & \code{STATIC\_PRESSURE\_RATIO}
\end{tabular}}
\begin{itemize}\small
  \item On the unsteady fluid-plot path, each parameter becomes one plot per vertex:
        \code{<PARAMETER><n>}, with \code{n} counting across entries.
        The probes table un-pivots these columns to one row per probe and step.
  \item On a steady run, \code{EXPORT\_PROBE\_POINTS} takes no variable list.
        Its columns are the solver's own.
  \item An entry with empty \code{parameters} emits nothing on either run type.
\end{itemize}
\small These six boundary-layer parameters are accepted on builds
26.122, 26.123 and 26.124. An unsupported parameter is refused by name and build:
\par\vspace{2mm}
{\ftstablesize
\begin{tabular}{@{}ll@{}}
\code{BL\_MOMENTUM\_THICKNESS} & \code{BL\_DISPLACEMENT\_THICKNESS} \\
\code{BL\_TOTAL\_THICKNESS} & \code{BL\_SHAPE\_FACTOR} \\
\code{BL\_SKIN\_FRICTION} & \code{BL\_TRANSITION\_MARKER}
\end{tabular}}
\par\vfill
{\tiny\color{SoftGray} Sources under \code{src/pyflightstream/}:
\code{cases/\_\_init\_\_.py}, \code{FLUID\_PLOT\_PARAMETERS}, \code{ProbesSpec};
\code{commands/unsteady\_solver.yaml}, \code{UNSTEADY\_SOLVER\_NEW\_FLUID\_PLOT};
\code{commands/probe\_points.yaml}, \code{EXPORT\_PROBE\_POINTS};
\code{cases/workflows.py}, \code{\_emit\_one\_probe\_table};
\code{post/products.py}, \code{write\_unsteady\_probes\_table}.\par}
\end{frame}

\begin{frame}{The probe source follows the run type}
\begin{itemize}\small
  \item On an unsteady row, every \code{[[probes]]} entry uses fluid plots:
        drawn lines and cited \code{points\_file} profiles alike.
  \item Each vertex and selected parameter shares one plot counter.
        \code{probes/<point>\_probes.csv} holds both forms together,
        one row per probe and time step, from the plots history.
  \item A cited profile is read at plan: a count, then \code{X,Y,Z,TYPE}
        CSV rows. Types 0 and 1 become fixed fluid-plot vertices with the
        entry's frame and scale. Invalid profiles are refused before solving.
  \item The unsteady path no longer imports or exports standard probe points.
        Its default output list has seven files instead of eight.
\end{itemize}
\end{frame}

\begin{frame}{Steady probes and older unsteady records}
\begin{itemize}\small
  \item Steady rows keep the probe-points source for both drawn and cited probes.
        A nonempty \code{parameters} list enables the entry; an empty list disables it.
  \item On a steady row the list does not filter exported variables:
        the solver exports a fixed set. Plan warns with the entry number and frame.
  \item On an unsteady row the list selects the sampled fluid-plot variables.
  \item An older cited profile exported only as an instant has no history.
        Re-post skips it with its reason in \code{products.json}, preserving
        the drawn-line histories that were recorded. A new run supplies
        the missing history.
\end{itemize}
\end{frame}

\begin{frame}[fragile]{What a point exports: \code{[exports]}}
\relax
{\ftstablesize
\begin{tabular}{@{}p{0.17\linewidth}p{0.13\linewidth}p{0.51\linewidth}p{0.098\linewidth}@{}}
\toprule
\textbf{Key} & \textbf{File suffix} & \textbf{Solver command} & \textbf{Run type} \\
\midrule
\code{simulation} & \code{.fsm} & \code{SAVEAS} & Both \\
\code{loads} & \code{.txt} & \code{EXPORT\_SOLVER\_ANALYSIS\_SPREADSHEET} & Both \\
\code{tecplot} & \code{.dat}, \code{.vtk} & \code{EXPORT\_SOLVER\_ANALYSIS\_VTK}; the package writes the \code{.dat} & Both \\
\code{sections} & \code{\_cp.txt} & \code{EXPORT\_ALL\_SURFACE\_SECTIONS} & Both \\
\code{sectional\_loads} & \code{\_sloads.txt} & \code{EXPORT\_SURFACE\_SECTIONAL\_LOADS} & Both \\
\code{probes} & \code{\_probes.txt} & \code{EXPORT\_PROBE\_POINTS} & Steady only \\
\code{plots} & \code{\_plots.txt} & \code{UNSTEADY\_SOLVER\_EXPORT\_PLOTS} & Unsteady only \\
\code{log} & \code{\_log.txt} & \code{EXPORT\_LOG} & Both \\
\bottomrule
\end{tabular}}
\begin{itemize}\small
  \item Keys are \code{true} or \code{false}. The original kinds default on for their run type.
        \code{vtk}, \code{csv} and \code{force\_distributions} default off; the
        steady plot kinds default on (next frame).
\end{itemize}
\end{frame}

\begin{frame}{What a point exports: \code{[exports]} (continued)}
\begin{itemize}\small
  \item \code{loads} is required to judge the run, and \code{simulation}
        is required too: every point leaves its final
        \code{.fsm}. Either one stated \code{false} is refused.
  \item \code{tecplot} exports every surface as VTK and the package writes the
        \code{.dat} from it, per panel, in the reference frame.
  \item \code{singularity\_strength = true}, a top-level pproc key, off by
        default, adds the nodal \code{Singularity\_strength} of a native
        auxiliary \code{*\_native\_tecplot.dat}, matched to the same step's
        VTK by geometry [4]. Off, the point makes no native export, its
        \code{.dat} states the strength \code{NOT\_CARRIED} (so does its
        \code{products.json} entry), and the point is complete without it.
        \pyfs{plan} says per row whether the strength is carried.
\end{itemize}
{\tiny\color{SoftGray} Sources under \code{src/pyflightstream/}:
\code{cases/\_\_init\_\_.py}, \code{EXPORT\_KINDS}, \code{default\_outputs},
\code{\_known\_export\_kinds}; \code{commands/solver\_export.yaml},
\code{EXPORT\_SOLVER\_ANALYSIS\_VTK}; \code{results/surface.py}.\par}
\end{frame}

\begin{frame}{Additional exports: plots, a volume section, force distributions}
\begin{itemize}\small
  \item \textbf{The solver's own plots.} \code{plot\_residuals},
        \code{plot\_loads} (on by default) and \code{plot\_sections\_cp} (on
        where sections are declared) under \code{[exports]}: the plotted series
        as TEXT, one row per iteration or one x/Cp pair per section.
        \textbf{Never a coefficient source}: on 26.124 the plotted pitching
        moment is $-0.897$ where the exported \code{CMy} is $-0.0247$ [3].
        An unsteady row saves the first two as well, once,
        after the march, and the section Cp too, where
        sections are declared.
        \code{false} turns each off.
  \item \textbf{A volume section.} \code{[volume\_section]}: one plane, a
        rectangle (\code{corners\_m}) or an annulus (\code{radii\_m}), in a
        \code{plane} of a \code{frame}, at \code{offset\_m}. \alert{It
        is sampled, not exported by the solver}: probes on a steady
        row, fluid plots on an unsteady or rotor row, written as
        a vertex cloud to \code{post/<matrix>/fields/<point>\_vsec.vtk} or
        \code{.dat} (\code{\_step\_<STEP>} per unsteady STEP). Older records
        keep their native export in the datapoint.
  \item \textbf{The force distribution.} \code{force\_distributions =
        true}: per-panel pressure and viscous force coefficients, once, at the
        end of the run. Off by default.
\end{itemize}
\end{frame}

\begin{frame}[fragile]{VTK and CSV surface exports}
\small \code{vtk\_variables} is a top-level pproc key, before any tables:
\begin{lstlisting}[style=ftsbase]
vtk_variables = ["X", "Y", "Z", "CP_FREESTREAM"]
[exports]
vtk = true
csv = true
\end{lstlisting}
\begin{itemize}\small
  \item Both kinds default off. VTK uses \code{EXPORT\_SOLVER\_ANALYSIS\_VTK}
        and \code{SET\_VTK\_EXPORT\_VARIABLES} when the list is present.
        An absent list requests all variables.
  \item CSV uses \code{EXPORT\_SOLVER\_ANALYSIS\_CSV}: \code{CP-FREESTREAM},
        \code{PASCALS}, all surfaces, solver reference frame. Both exclude the wake.
  \item Commands and variables are checked against the selected build's database.
        Unsupported requests are refused at plan, naming the build.
\end{itemize}
\end{frame}

\begin{frame}[fragile]{The surface averaging window: the package averages}
\begin{lstlisting}[style=ftsbase]
[time_averaging]
last_revs = 1.5 # OR last_iters = 54; exactly one, positive
\end{lstlisting}
\begin{itemize}\small
  \item \alert{The package averages; the solver does not.} The run
        exports the surface at every step of the window, as
        \code{EXPORT\_UNSTEADY\_AFTER\_ITER: <first step>} would, and the post
        averages those exports into \code{surfaces/<point>\_time\_average.dat}:
        the same panel across the steps, every step the same weight, every
        variable. \code{SOLVER\_TIME\_AVERAGING} is never emitted: it hung
        26.124, killed at 240.5 seconds with nothing written [3], and 26.123
        does not carry it.
\end{itemize}
\end{frame}

\begin{frame}{The surface window: time steps, verified on 26.122}
\begin{itemize}\small
  \item Bounds are inclusive, 1-based time steps ending at the run's last step.
        \code{last\_iters} counts steps; \code{last\_revs} uses the same rotor
        clock and rounding as \code{LAST\_REVS\_AVG}. A long window clips at step 1.
  \item Verified on 26.122 [3]: the solver's own average over steps
        7 to 12 is the uniform mean of the run's per-step surfaces over those
        steps, to 1e-13 in Cp, Vx and speed, which is the package's mean. The
        solver keeps CF at its last instant; the package averages it.
  \item Steps that do not share one topology refuse the average; a step of the
        window that was not exported (a clock-stopped run, a continuation) skips
        it. Both by name, never a partial average.
  \item Needs 26.122 or later and \code{tecplot} on. The export header's
        inner-iteration count is not this clock.
\end{itemize}
\end{frame}

\begin{frame}{Surface export provenance and per-step exports}
\begin{itemize}\small
  \item \code{products.json} states the averaged surface \code{kind: average},
        with the recorded window, the steps averaged and each per-step VTK read
        with its sha256; the per-step exports stay \code{kind: instant}.
  \item The window is read from the run record (\code{surface\_average\_window}),
        never from a pproc edited since; changing it requires a new run.
  \item Every \code{.dat} states its source in \code{DATASETAUXDATA}: the VTK,
        its sha256 and a \code{TRANSLATION} record (cell-centred, reference frame).
  \item \code{EXPORT\_UNSTEADY\_AFTER\_REV} or \code{EXPORT\_UNSTEADY\_AFTER\_ITER}
        exports the VTK per step, CSV where asked; each step's \code{.dat} is
        written from its VTK.
  \item An older record keeps the solver's window and its
        meaning. The matrix's plots-reduction window remains a separate choice.
  \item A storage recipe's \code{[[prune\_step\_exports]]} keeps only each
        point's last step; a post that needs a deleted step refuses, naming it
        (Guide 2).
\end{itemize}
\end{frame}

\begin{frame}{Surface formats without a pproc export key}
\begin{itemize}\small
  \item \code{EXPORT\_SOLVER\_ANALYSIS\_PLOAD\_BDF} still has no
        \code{[exports]} key. \code{EXPORT\_SOLVER\_ANALYSIS\_FORCE\_DISTRIBUTIONS}
        has one, \code{force\_distributions}.
  \item Its parameter or payload block cannot be supplied by a setup's
        one-line \code{[[raw]]} entry.
  \item \code{EXPORT\_BL\_VELOCITY\_PROFILE} holds an unattended script
        [3]: a raw line is refused at plan on 26.124.
\end{itemize}
\end{frame}
